\section{Introdução}

APIs Web são o principal meio de integração entre sistemas modernos.
Duas abordagens dominam o cenário atual: REST (Representational State Transfer)
e GraphQL, uma linguagem de consulta desenvolvida pelo Facebook em 2015.
Enquanto REST expõe múltiplos endpoints com estruturas fixas, GraphQL permite
que o cliente especifique exatamente os campos desejados em uma única requisição.

Este experimento investiga se GraphQL oferece vantagens mensuráveis sobre REST
em dois aspectos críticos para aplicações web: velocidade de resposta e
eficiência no uso da rede.

\subsection{Perguntas de Pesquisa}

\textbf{RQ1.} Respostas às consultas GraphQL são mais rápidas que respostas às
consultas REST?

\textbf{RQ2.} Respostas às consultas GraphQL têm tamanho menor que respostas
às consultas REST?

\subsection{Hipóteses}

Para RQ1:
\[
H_0: \mu_{\text{GraphQL}} \ge \mu_{\text{REST}} \quad
H_1: \mu_{\text{GraphQL}} < \mu_{\text{REST}}
\]

Para RQ2:
\[
H_0: \mu_{\text{GraphQL}} \ge \mu_{\text{REST}} \quad
H_1: \mu_{\text{GraphQL}} < \mu_{\text{REST}}
\]

Nível de significância: $\alpha = 0,05$.
